\clearpage
\begin{explanationbox}
\textbf{先看一个真点：}$y^2=4x$ 上的 $P(4,4)$，焦点为 $F(1,0)$。
从 $P$ 向准线作垂线，垂足为 $Q$；与 $y$ 轴交于 $H$。
\begin{center}% Exact representative: p=2, P=(4,4), F=(1,0), PF=PQ=5.
\begin{tikzpicture}[>=stealth,font=\small]
\draw[rounded corners=3pt,gray!45] (0,0) rectangle (14.2,10.4);
\node at (7.1,9.85) {斜线换水平线：$PF=PQ=QH+HP$};
\begin{scope}[shift={(3.3,3.4)},scale=0.88]
  \draw[->,gray] (-1.8,0)--(5.9,0) node[right] {$x$};
  \draw[->,gray] (0,-2.6)--(0,5.5) node[above] {$y$};
  \draw[red!75!black,thick] (-1,-2.4)--(-1,5.3);
  \draw[teal!75!black,thick,domain=-2.3:4.6,samples=100] plot ({\x*\x/4},\x);
  \draw[blue!75!black,very thick] (1,0)--(4,4) node[pos=.48,right=3pt] {$PF$};
  \draw[violet,very thick] (-1,4)--(0,4);
  \draw[orange!85!black,very thick] (0,4)--(4,4);
  \draw[gray,dashed] (4,0)--(4,4);
  \draw[red!75!black] (-1,3.7)--(-.7,3.7)--(-.7,4);
  \draw[<->,violet] (-1,4.75)--(0,4.75) node[midway,above] {$p/2$};
  \draw[<->,orange!85!black] (0,4.75)--(4,4.75) node[midway,above] {$x_0$};
  \fill (-1,4) circle(2pt) node[left=3pt] {$Q$};
  \fill (0,4) circle(2pt) node[below left=2pt] {$H$};
  \fill (4,4) circle(2pt) node[right=4pt] {$P(x_0,y_0)$};
  \fill (0,0) circle(2pt) node[below left=3pt] {$O$};
  \fill (1,0) circle(2pt) node[below=4pt] {$F$};
  \node[below=4pt] at (4,0) {$x_0$};
  \node[left,align=center] at (-1,-1.5) {准线 $l$\\$x=-p/2$};
\end{scope}
\node[align=left,anchor=west] at (10.2,6.6) {$P$ 在抛物线上\\[6pt]$\color{blue!75!black}PF=PQ$\\[12pt]$QH=p/2$\\[6pt]$HP=x_0$};
\node at (7.1,1.15) {$\displaystyle PF=\textcolor{violet}{\frac p2}+\textcolor{orange!85!black}{x_0}$};
\node at (7.1,.5) {图取 $y^2=4x$、$P(4,4)$：两段 $1+4$，焦半径就是 $5$。};
\end{tikzpicture}
\end{center}
\textbf{只看两件事：}点在抛物线上，所以蓝色焦半径 $PF=PQ$；水平的 $PQ$ 又恰好分成紫段 $p/2$ 和橙段 $x_0$。
\[
  \boxed{PF=PQ=QH+HP=\frac p2+x_0}.
\]
图取 $p=2$，但分段关系对任意 $p>0$、任意曲线上点都成立（顶点处 $H=P$）。
\textbf{记忆动作：}遇到焦半径，先找准线，再量沿开口方向的距离。
\end{explanationbox}
\clearpage
\begin{explanationbox}
\textbf{公式里为什么加 $p/2$？}顶点在焦点与准线的正中间。
\begin{center}\begin{tikzpicture}[scale=0.95,>=stealth]
\draw[blue,thick,domain=-3:3.5,samples=90] plot ({\x*\x/4},\x);\draw[red,thick] (-1,-3.2)--(-1,3.5);\draw[->,gray] (-1,0)--(2,0);\fill (-1,0) circle (1.7pt) node[below left] {$D$};\fill (0,0) circle (1.7pt) node[below left=3pt] {$O$};\fill (1,0) circle (1.7pt) node[below right=3pt] {$F$};\draw[gray,dotted] (0,0)--(0,-3.8);\draw[gray,dotted] (1,0)--(1,-3.8);\draw[<->,orange,thick] (0,-3.8)--(1,-3.8) node[midway,below] {$OF=p/2=1$};\draw[<->,teal,thick] (1,-2)--(1,2);\draw[gray] (1,2)--(3.5,2.1);\node[right] at (3.7,2.1) {通径 $2p=4$};
\end{tikzpicture}
\end{center}
\textbf{参数不混淆：}$DF=p$，$DO=OF=p/2$；通径 $AB=2p$，半通径 $AF=FB=p$。
因此公式中的常数是 $OF$，不是半通径。

\medskip
\textbf{边界也能看见：}$PF=x_0+p/2\ge p/2$，且只有 $P=O$ 时取等号。
已知 $PF=r$ 时可逆求 $x_0=r-p/2$，须有 $r\ge p/2$。
若写成 $y^2=4ax$，则 $a=p/2$，公式相应为 $PF=x_0+a$。
\end{explanationbox}
